\section{Related work}

\paragraph{Pocket-conditioned generation.} Structure-based generative models place atoms in a
measured site and several are explicitly interaction-aware: Pocket2Mol~\citep{peng2022pocket2mol},
TargetDiff~\citep{guan2023targetdiff}, DiffSBDD~\citep{schneuing2024diffsbdd},
DrugFlow~\citep{schneuing2025drugflow}, ResGen~\citep{resgen2023}, PocketFlow~\citep{jiang2024pocketflow}
and Lingo3DMol~\citep{feng2024lingo3dmol}, the last predicting non-covalent interaction sites
directly. They do not state a design hypothesis in words, and their evaluation rests on
distributional and docking-derived scores whose adequacy has been
criticised~\citep{cieplinski2023dockwell,harris2023posecheck,buttenschoen2024posebusters}. Recent
work does prompt an LLM with pockets and required residue contacts~\citep{macdougall2026spatial}, so
we do not claim that reasoning LLMs cannot see a pocket; we claim that no prior system reports
whether its stated reasons are checkable against the ligand's own contact table.

\paragraph{Reasoning over molecular optimisation.} Chemistry LLMs and reasoning models now justify
modifications in natural language~\citep{zhao2024chemdfm,narayanan2025ether0,cavanagh2024smileyllama,
fang2023molinstructions,zhang2024chemllm}. Closest to us, DESRO~\citep{liu2026desro} recovers
reasoning steps from outcome data, DrugR~\citep{liu2026drugr} builds rationales backwards from
computed ADMET outcomes, and Logos~\citep{wen2026logos} emits a reason block then a SMILES. All
three generate reason and molecule inside one model and one trace, so the reason cannot be
substituted and the intervention we rely on is unavailable; none conditions on pocket structure.
MolecularCanvas~\citep{dong2026molecularcanvas} separates an LLM plan from an RDKit executor but is
an interactive tool without held-out oracle or mismatch measurements. Earlier work factorised
rationale-to-molecule where the rationale is a \emph{substructure} rather than a
sentence~\citep{jin2020rationale}. Our contribution against this group is the separation into two
independently trained models with a natural-language interface, the structural grounding of that
interface, and the control arms that quantify it.

\paragraph{Instruction-faithfulness controls.} Pocket-shuffled controls are standard in
structure-based design, and UTGDiff~\citep{utgdiff2024} argues that distributional metrics cannot
detect cross-matched instruction/molecule pairs. Our three-arm decomposition adds the finding that a
mismatched instruction falls \emph{below} the no-instruction arm, and that the no-instruction arm is
degenerate because it echoes the input.

\paragraph{Lead optimisation from measured pairs.} Delta and pairwise models exploit analogue pairs
directly~\citep{fralish2023deepdelta,feng2024actfound}, matched molecular series have long been used
to pick the next compound~\citep{wawer2014whatnext}, and REINVENT~4~\citep{loeffler2024reinvent4}
provides the analogue-generation baseline at matched budget. These supply molecules and potency
estimates, not hypotheses.

\paragraph{Literature mining.} LLM extraction from chemistry papers targets reactions, procedures
and entities~\citep{chemminer2024,materials2025ie}. We are not aware of prior work extracting the
stated design rationale of a medicinal-chemistry modification from paper text.

\paragraph{Cofolding as design input.} We take contact tables from cofolded complexes where no
crystal structure exists~\citep{wohlwend2024boltz1,passaro2025boltz2,abramson2024af3,chai2024chai1},
and we adopt the caveats from prospective cofolding
evaluations~\citep{cofold_prospective_2025} rather than treating a predicted pose as a measurement.
